% Standalone final source is assembled from reviewed working fragments.
\documentclass[11pt,a4paper]{article}
\usepackage[T1]{fontenc}
\usepackage{lmodern,microtype,amsmath,amssymb,amsthm,mathtools}
\usepackage[margin=24mm,headheight=15pt]{geometry}
\usepackage{booktabs,array,xcolor,fancyhdr,titlesec,fancyvrb,tikz}
\usepackage[colorlinks=true,linkcolor=teal!65!black,citecolor=teal!65!black,urlcolor=teal!65!black]{hyperref}
\definecolor{ink}{HTML}{15364A}
\definecolor{accent}{HTML}{007C78}
\titleformat{\section}{\Large\sffamily\bfseries\color{ink}}{\thesection}{.55em}{}
\titleformat{\subsection}{\large\sffamily\bfseries\color{accent}}{\thesubsection}{.55em}{}
\newtheorem{theorem}{Theorem}[section]
\newtheorem{lemma}[theorem]{Lemma}
\newtheorem{proposition}[theorem]{Proposition}
\newtheorem{corollary}[theorem]{Corollary}
\theoremstyle{remark}\newtheorem{remark}[theorem]{Remark}
\DeclareMathOperator{\conv}{conv}
\DeclareMathOperator{\vol}{vol}
\DeclareMathOperator{\Vol}{Vol}
\DeclareMathOperator{\inter}{int}
\DeclareMathOperator{\relint}{relint}
\DeclareMathOperator{\diam}{diam}
\newcommand{\R}{\mathbb R}\newcommand{\Z}{\mathbb Z}
\newcommand{\fr}[1]{\left\{#1\right\}}\newcommand{\one}{\mathbf1}
\pagestyle{fancy}\fancyhf{}
\fancyhead[L]{\small\sffamily Exact lattice covering transfers}
\fancyhead[R]{\small\sffamily Revised research manuscript}
\fancyfoot[L]{\footnotesize\sffamily Prepared for Keio | 20 November 2026}
\fancyfoot[R]{\small\thepage}
\setlength{\parindent}{1.1em}\setlength{\parskip}{3pt}
\setlength{\emergencystretch}{2em}\raggedbottom
\hypersetup{pdftitle={Exact Lattice Covering Transfers: Covering Minima, Partitions, and Arithmetic Families},pdfauthor={Research manuscript prepared for Keio}}

\begin{document}
\noindent{\sffamily\small\bfseries\color{accent} GEOMETRY OF NUMBERS | TRANSFER FROM VMRC [200]}\par\medskip
\noindent{\sffamily\bfseries\fontsize{23}{27}\selectfont\color{ink}\raggedright
Exact Lattice Covering Transfers:\\
Covering Minima, Partitions,\\
and Arithmetic Families\par}
\medskip
\noindent{\large\sffamily All real parameters, every ambient dimension,\par
and exact finite formulas beyond an upper bound}\par
\medskip
\noindent Research manuscript prepared for Keio\hfill 20 November 2026

\begin{abstract}
We develop exact consequences of the weighted fractional-part method
in VMRC [200]. First, for every integer $n\ge3$ and every real $r\ge1$,
we prove $\mu_2(S_n(1,1,r,\ldots,r))=\tfrac12+\tfrac1{r+1}$.
This gives an exact counterexample family to the weighted covering-minima
conjecture of Codenotti--Santos--Schymura, uniformly in ambient dimension,
and closes the nonintegral-parameter gap in our first transfer note.
The main new upper-bound ingredient is a planar theorem allowing all
but at most two lattice columns to be radially expanded.
Second, for positive relations $(k,1,\ldots,1)$ we express the entire
covering radius as an explicit finite optimization over ordered
partitions. For $k=2$ the deficit from the surface-area upper bound
is exactly proportional to minimum subset-sum discrepancy. Equal
singleton weights give a closed formula involving $\gcd(d,k)$.
We obtain an exact all-dimensional family satisfying the numerical
interior-point conjecture, prove closure of that inequality under
direct sums with interior origins, and show globally that a non-hollow
polytope with covering radius $d/2$ has exactly one interior lattice
point. We also give a finite linear characterization of weighted
equality and retain explicit counterexamples to the proposed equality
classifications. The full numerical interior-point conjecture remains
unresolved by this manuscript.
\end{abstract}

\section{Questions, main results, and scope}
For a convex body $K\subset\R^d$ and a full-rank lattice $\Lambda$,
the covering radius is
\[
 \mu(K,\Lambda)=\inf\{\rho\ge0:\rho K+\Lambda=\R^d\}.
\]
The $i$th covering minimum is the least dilation whose lattice
arrangement meets every affine subspace of codimension $i$.
We use its standard quotient characterization \cite[Section 1]{CSS}:
\begin{equation}\label{intro:quotient}
 \mu_i(K,\Lambda)=\max_A\mu(AK,A\Lambda),
\end{equation}
where $A$ ranges over rank-$i$ linear quotients with rational kernel
relative to $\Lambda$. The image lattice is always the actual
projected lattice $A\Lambda$. Covering radii are unchanged by
translations of the body and simultaneous invertible linear changes
of body and lattice. They scale inversely under positive dilations.

For positive weights define
\[
 S_n(\omega)=\conv\{-\omega_0\one_n,
                      \omega_1e_1,\ldots,\omega_ne_n\}.
\]
The weighted conjecture \cite[Conjecture 5.3]{CSS}, also explicitly
repeated in January 2026 \cite[Conjecture 5.3]{Krivokuca}, asserts,
for ordered weights $\omega_0\le\cdots\le\omega_n$,
\begin{equation}\label{intro:weightedconj}
 \mu_i(S_n(\omega))=
 \frac{\displaystyle\sum_{0\le j<\ell\le i}(\omega_j\omega_\ell)^{-1}}
      {\displaystyle\sum_{j=0}^i\omega_j^{-1}}.
\end{equation}
Our exact answer for the family with two smaller weights is
\begin{equation}\label{intro:mainmin}
 \boxed{\mu_2(S_n(a,a,b,\ldots,b))
       =\frac1{2a}+\frac1{a+b}\quad(n\ge3,\ 0<a\le b).}
\end{equation}
It strictly exceeds \eqref{intro:weightedconj} when $a<b$.
In particular the smallest integral example is
$\mu_2(S_3(1,1,2,2))=5/6>4/5$.
The cases $i=1$ and $i=n$ of the conjecture are prior theorems;
dimension three is therefore the smallest possible counterexample
dimension. The present revision proves \eqref{intro:mainmin} for
\emph{all} real parameters and \emph{all} $n\ge3$.

The second principal result concerns a different simplex family,
whose positive relation has one coefficient $k$ and all others one.
Its covering radius is computed exactly by an ordered-partition
cost, with no phase optimization left over. A two-group specialization
is a direct geometric realization of minimum partition discrepancy.
When the singly occurring weights are equal, the formula reduces to
a greatest-common-divisor expression. These statements concern exact
values and equality, extending the phase-uniform estimate of [200].

For interior lattice points, CSS Conjecture D proposes
\begin{equation}\label{intro:D}
 \mu(P)\le\frac{d-1}{2}+\frac1{h+1},\qquad
 h=|\inter P\cap\Lambda|\ge1.
\end{equation}
We prove it for an explicit two-parameter simplex family in every
dimension and for direct sums of bodies already satisfying it.
We also prove the weaker but fully global implication
\[
 |\inter P\cap\Lambda|\ge2\quad\Longrightarrow\quad\mu(P)<d/2.
\]
The precise quantitative bound \eqref{intro:D} for arbitrary lattice
polytopes is a remaining question. The equality-classification
counterexamples from the first transfer note are included with proofs.

\section{Exact transfer tools from fractional-part flows}
\label{sec:tools}
For $0\in\inter P$, write
$g_P(x)=\inf\{\rho\ge0:x\in\rho P\}$ for its gauge.
The following lemmas reproduce the needed machinery of VMRC [200]
\cite{VMRC200}, so the applications can be checked within this text.

\begin{lemma}[Exact weighted gauge]\label{tool:gauge}
Let $u_0,\ldots,u_d$ generate the lattice $\Lambda$, and let their
column matrix $U$ have rank $d$ and kernel $\R q$, where
$q\in\Z_{>0}^{d+1}$. For positive $c_j$ set
$P=\conv(u_0/c_0,\ldots,u_d/c_d)$. Then, for every $x=Ut$,
\begin{equation}\label{tool:gaugeeq}
 \min_{z\in\Lambda}g_P(x-z)
 =\min_{s\in[0,1)}\sum_{j=0}^dc_j\fr{t_j+q_js}.
\end{equation}
Consequently the covering radius is the maximum of the right side
over all phases $t$.
\end{lemma}
\begin{proof}
The positive relation puts zero in the interior of $P$.
Every lattice vector is $Un$ for $n\in\Z^{d+1}$, and every
coefficient representation of $x-Un$ is $t-n+sq$.
Its gauge is the infimum of $\sum_jc_j(t_j-n_j+sq_j)$ over these
representations with nonnegative coefficients. For fixed $s$,
the largest allowed $n_j$ is $\lfloor t_j+sq_j\rfloor$, which
minimizes the positive cost. Interchanging the infima over $s,n$
therefore gives the stated fractional parts. The result is
one-periodic and has positive slope between finitely many downward
jumps, so the minimum is attained. The lattice gauge distance is
continuous on the compact quotient torus; its maximum is the
covering radius and is attained as well.
\end{proof}

\begin{lemma}[Jump-square bound and equality]\label{tool:jump}
For positive integers $q_j$, positive real $c_j$, and arbitrary
phases $t_j$, put
\[
 F(s)=\sum_jc_j\fr{t_j+q_js},\qquad C=\sum_jq_jc_j.
\]
Then
\begin{equation}\label{tool:bound}
 \min_sF(s)\le B(q,c):=
 \frac12\sum_jc_j-\frac{\sum_jq_jc_j^2}{2C}.
\end{equation}
Equality holds exactly when all coordinate crossings are distinct
and all post-jump values equal $B(q,c)$.
\end{lemma}
\begin{proof}
Group simultaneous crossings into events $e$ with total jump $J_e$.
Between events $F'=C$, while
$\int_0^1F=\tfrac12\sum_jc_j$ and $\sum_eJ_e=C$.
The change of $F^2$ at $e$ is $-2J_eF(e)-J_e^2$, using its
right-continuous value. Periodicity thus gives
\[
 \frac1C\sum_eJ_eF(e)
 =\frac12\sum_jc_j-\frac1{2C}\sum_eJ_e^2
 \le B(q,c).
\]
The first side is a convex combination of post-jump values.
The inequality follows from positivity, because grouped squared
jumps are at least the sum of the individual squares.
Equality requires all values in the convex combination to be the
minimum and every group to have one crossing; these conditions
also suffice.
\end{proof}

\begin{remark}[Individual jumps at coincident times]\label{tool:individual}
For exact deficit calculations, simultaneous crossings may instead
be recorded in any specified sequential order, with zero time
between them. Let $w_e$ be the individual jump size and let $H_e$
be its formal intermediate post-jump value. These values are all
at least the actual final post-jump value. Telescoping the individual
changes of $F^2$ remains exact, and for $M=\min F$ it yields
\begin{equation}\label{tool:deficit}
 M=B(q,c)-\frac1C\sum_e w_e(H_e-M),\qquad H_e-M\ge0.
\end{equation}
This convention accounts for cross terms of coincident jumps; it
does not identify intermediate values with the right-continuous
value of $F$ at that time.
\end{remark}

\begin{corollary}[Upper bound without a lattice-generation assumption]
\label{tool:upper}
If integer columns $u_0,\ldots,u_d$ span $\R^d$ and have a positive
integer relation $q$, then for $P=\conv(u_j/c_j)$,
\[
 \mu(P,\Lambda)\le B(q,c),
\]
even when the columns are nonprimitive or generate a proper
sublattice of $\Lambda$.
\end{corollary}
\begin{proof}
For $x=\sum_jt_ju_j$, the vector
$z(s)=\sum_j\lfloor t_j+sq_j\rfloor u_j$ lies in $\Lambda$ and
\[
 x-z(s)=\sum_jc_j\fr{t_j+sq_j}(u_j/c_j).
\]
The coefficients are nonnegative. At a minimizing event their
sum is at most $B(q,c)$, so the residual lies in $B(q,c)P$.
This works for every real $x$. Only the exact reverse direction
in Lemma~\ref{tool:gauge} needs lattice generation.
\end{proof}

\begin{lemma}[Section and quotient]\label{tool:section}
Let $0\in\inter K$ and let $L$ be a rational linear subspace.
For the quotient map $\pi$,
\[
 \mu(K,\Lambda)\le
 \mu(K\cap L,\Lambda\cap L)+\mu(\pi K,\pi\Lambda).
\]
\end{lemma}
\begin{proof}
Take numbers $a,b$ strictly larger than the section and quotient
radii. For any $x$, quotient covering supplies $y\in bK$ and
$z\in\Lambda$ with $x-y-z\in L$. Section covering writes this
residual as $w+v$, where $w\in\Lambda\cap L$ and $v\in a(K\cap L)$.
Thus $x\in z+w+(a+b)K$ by convexity. Let $a,b$ decrease to their
radii. This is the standard inequality \cite[Lemma 2.1]{CSS}.
\end{proof}

\begin{proposition}[The non-hollow bound from the transfer]
\label{tool:nonhollow}
Every lattice $d$-polytope with an interior lattice point has
covering radius at most $d/2$.
\end{proposition}
\begin{proof}
For a lattice simplex with zero in its interior, apply
Corollary~\ref{tool:upper} to its actual vertices and all $c_j=1$.
The positive integer relation exists because the vertex matrix is
integral and has a positive one-dimensional kernel. The resulting
bound is $(d+1)/2-1/2=d/2$.

For a polytope, translate an interior lattice point to zero and
choose an inclusion-minimal set of vertices whose convex hull
contains zero. It is a simplex $S$ with zero in its relative
interior. Let $r=\dim S$ and $L=\operatorname{span}S$. If $r=d$,
monotonicity gives the simplex bound. Otherwise $\pi P$ is a
non-hollow lattice $(d-r)$-polytope, and induction together with
Lemma~\ref{tool:section} gives
\[
 \mu(P)\le\mu(P\cap L)+\mu(\pi P)
 \le\mu(S)+(d-r)/2\le d/2.
\]
The one-dimensional base case is the interval statement. This
reproduces the required consequence of [200] with the actual
ambient lattice throughout.
\end{proof}



\section{A planar stretching theorem}
\label{rr:sec:stretch}

The weighted fractional-part bound also controls a larger family of
planar bodies.  This will remove the integer-parameter restriction in
the first counterexample and determine the second covering minimum
in every ambient dimension.

For $r\ge1$, write
\begin{equation}\label{rr:eq:H}
 H_r=\frac12+\frac1{r+1}.
\end{equation}

\begin{lemma}[A two-frequency bound]\label{rr:lem:twofrequency}
For positive integers $p,q$ and arbitrary real $t_0,t_1$,
\[
 \min_{s\in\R}\bigl(\fr{t_0+ps}+\fr{t_1+qs}\bigr)
 \le\frac{\gcd(p,q)}{p+q}.
\]
\end{lemma}
\begin{proof}
Put $g=\gcd(p,q)$.  The integers $v_0=q/g$ and $v_1=-p/g$
generate $\Z$ and have positive primitive relation $(p/g,q/g)$.
By the exact gauge formula, Lemma~\ref{tool:gauge}, the displayed
minimum is the distance of
$x=(q/g)t_0-(p/g)t_1$ from $\Z$ in the gauge of
$[-p/g,q/g]$.  Replacing the flow parameter by $gs$ does not change
the minimum.  The interval has lattice length $(p+q)/g$, hence
covering radius $g/(p+q)$.
\end{proof}

We will use the following consequence of Lemma~\ref{tool:jump}.
If lattice vectors $v_0,v_1,v_2$ span the plane and have a positive
integral relation $\sum q_jv_j=0$, then, for $c_j>0$,
\begin{equation}\label{rr:eq:trianglebound}
 \mu\bigl(\conv(v_0/c_0,v_1/c_1,v_2/c_2),\Lambda\bigr)
 \le B(q,c):=
 \frac12\sum_jc_j-
 \frac{\sum_jq_jc_j^2}{2\sum_jq_jc_j}.
\end{equation}
Neither primitivity of the individual vectors nor generation of the
ambient lattice is required for this upper bound.  Indeed, write
$x=\sum t_jv_j$ and subtract
$\sum\lfloor t_j+sq_j\rfloor v_j\in\Lambda$.  The remaining vector
is a nonnegative combination of the scaled vertices with coefficient
sum $\sum c_j\fr{t_j+sq_j}$.  Lemma~\ref{tool:jump} bounds its
minimum by $B(q,c)$.  The stronger lattice-generation hypothesis is
needed for the exact gauge identity, not for this rounding argument.

\begin{lemma}[Two fixed vertices of a lattice triangle]
\label{rr:lem:triangle}
Let $v_0,v_1,v_2$ belong to a planar lattice $\Lambda$, span the
plane, and satisfy $0\in\inter\conv(v_0,v_1,v_2)$.  For every real
$r\ge1$,
\[
 \mu\bigl(\conv(v_0,v_1,rv_2),\Lambda\bigr)\le H_r.
\]
\end{lemma}
\begin{proof}
Let $(p,q,\ell)$ be the primitive positive integral relation among
$v_0,v_1,v_2$, put $s=p+q$, and set $b=1/r$.  Thus
$0<b\le1$ and
\[
 H_r=\frac{1+3b}{2(1+b)}.
\]

\paragraph{The case $\ell\le s$.}
Use \eqref{rr:eq:trianglebound} with $c=(1,1,b)$.  It gives
\[
 B=\frac12\left(2+b-\frac{s+b^2\ell}{s+b\ell}\right),
 \qquad
 H_r-B=\frac{b(1-b)(s-\ell)}{2(1+b)(s+b\ell)}\ge0.
\]

\paragraph{The case $\ell>s$ and $1\le r\le2$.}
The unscaled lattice triangle $T=\conv(v_0,v_1,v_2)$ has two
distinct interior lattice points, $0$ and $-v_2$, because
\[
 -v_2=\frac p\ell v_0+\frac q\ell v_1+
       \left(1-\frac{s}{\ell}\right)0
\]
has positive coefficient on the interior point $0$.
The triangle $\conv(v_0,v_1,rv_2)$ contains $T$.  The established
planar covering theorem \cite[Theorem 6.3]{CSS} therefore yields
\[
 \mu\bigl(\conv(v_0,v_1,rv_2),\Lambda\bigr)
 \le\mu(T,\Lambda)\le\frac56\le H_r.
\]
Here the cited theorem says that a lattice polygon with $k\ge2$
interior lattice points has covering radius at most
$1/2+1/(k+1)$.

It remains to consider $\ell>s$ and $r\ge2$, so $b\le1/2$.

\paragraph{Unequal short-coordinate frequencies: $p\ne q$.}
By Lemma~\ref{rr:lem:twofrequency}, for every phase vector the
first two fractional parts can be made to sum to at most
$\gcd(p,q)/(p+q)\le1/3$.  The third term has weight $b$ and is at
most $b$.  Thus the rounding argument gives an upper bound
$1/3+b$, and
\[
 H_r-\left(\frac13+b\right)
 =\frac{(1+3b)(1-2b)}{6(1+b)}\ge0.
\]

\paragraph{Equal frequencies $p=q=k\ge2$.}
Primitivity gives $\gcd(k,\ell)=1$.  Choose $s_0$ so that the
two short-coordinate fractional parts sum to at most $1/2$.
Replacing $s_0$ by $s_0+j/k$, $0\le j<k$, preserves both of these
fractional parts.  The third fractional part visits a complete
equally spaced $k$-point grid on the circle, so one of its values
is at most $1/k$.  Hence the weighted sum is at most
\[
 \frac12+\frac bk\le\frac{1+b}{2}\le H_r.
\]

\paragraph{The remaining case $p=q=1$.}
Now $\ell\ge3$.  Pass to the sublattice
$\Lambda'=\Z v_0+\Z v_1+\Z v_2$, which can only increase the
covering radius.  Since $v_1=-v_0-\ell v_2$, the pair $v_2,v_0$
is a basis of $\Lambda'$.  In this basis the triangle becomes
\[
 R_{\ell,r}=\conv\{(-\ell,-1),(0,1),(r,0)\}.
\]
If $\ell=2k$, $k\ge2$, the affine unimodular map
$(x,y)\mapsto(y,x-ky+k)$ sends this triangle to
\[
 \conv\{(-1,0),(1,0),(0,r+k)\}.
\]
Its covering radius is $1/2+1/(r+k)\le H_r$; the elementary
radius calculation is supplied in Lemma~\ref{rr:lem:QL} below.

If $\ell=2k+1$, $k\ge1$, the affine unimodular map
$(x,y)\mapsto(x-ky+k,y)$ sends it to
\[
 \conv\{(-1,-1),(0,1),(r+k,0)\}=S_2(1,r+k,1).
\]
The known weighted-simplex radius formula
\cite[Theorem 1.4]{CSS} gives
\[
 \mu(R_{\ell,r},\Z^2)
 =\frac{r+k+2}{2(r+k)+1}
 \le\frac{r+3}{2r+3}
 <\frac{r+3}{2r+2}=H_r.
\]
The first inequality uses the fact that $(z+2)/(2z+1)$ decreases
for $z>0$.  This exhausts the possible relations.
\end{proof}

\begin{theorem}[Planar stretching with at most two fixed columns]
\label{rr:thm:stretch}
Let $u_1,\ldots,u_N$ be vectors of a planar lattice $\Lambda$ such
that $0\in\inter\conv(u_1,\ldots,u_N)$.  Keep at most two of these
vectors fixed and multiply every other vector by a real number
$r\ge1$.  If $P_r$ is their resulting convex hull, then
\begin{equation}\label{rr:eq:stretch}
 \mu(P_r,\Lambda)\le\frac12+\frac1{r+1}.
\end{equation}
The vectors need not be primitive, distinct, or generators of
$\Lambda$.
\end{theorem}
\begin{proof}
Positive radial rescaling preserves origin-interiority.  Zero
columns can be deleted, since the origin is already in the interior
of the convex hull of the nonzero columns.

Choose a positive relation of minimal support among the remaining
columns.  Its support has size at most three: in rank two, a
positive dependence on at least four columns can be perturbed
within the kernel until a coefficient vanishes.  Minimality
therefore leaves either a triangle containing zero in its interior
or a pair of opposite vectors.

For a triangle, there are at most two fixed columns.  With two,
apply Lemma~\ref{rr:lem:triangle}.  With one, put $b=1/r$ and
use \eqref{rr:eq:trianglebound} with $c=(1,b,b)$: the weighted
mean subtracted there is at least $b$, so the bound is at most
$(1+b)/2\le H_r$.  With no fixed column, $c=(b,b,b)$ and the
bound is $b\le H_r$.  Since the corresponding scaled triangle is
contained in $P_r$, each case proves \eqref{rr:eq:stretch}.

For an opposite pair, write the two vectors as $Aw$ and $-Bw$,
where $w$ is primitive in $\Lambda$ and $A,B$ are positive
integers.  Set $W=\R w$ and let $\pi$ be quotient by $W$.
In a primitive coordinate for $\pi\Lambda$, every projected
unscaled column is an integer.  The interval $\pi P_r$ has both
positive and negative endpoints because zero is interior to
$P_r$.  Since all multipliers are at least one, its lattice
length is at least two.  If all surviving projected columns are
scaled by $r$, its lattice length is at least $2r$.

Apply section/projection subadditivity,
Lemma~\ref{tool:section}.  If both selected vectors are scaled,
the section contains an interval of lattice length
$r(A+B)\ge2r$, so
\[
 \mu(P_r,\Lambda)\le\frac1{2r}+\frac12\le H_r.
\]
If exactly one is fixed, the section length is at least
$A+rB\ge1+r$, after naming $A$ as the fixed coefficient, and
\[
 \mu(P_r,\Lambda)\le\frac1{1+r}+\frac12=H_r.
\]
Finally, if both are fixed, the section length is at least two.
There are no other fixed columns, so the quotient interval has
length at least $2r$, giving
\[
 \mu(P_r,\Lambda)\le\frac12+\frac1{2r}\le H_r.
\]
These cases exhaust every minimal positive relation.
\end{proof}

\section{Exact second covering minima in every dimension}
\label{rr:sec:exact}

\begin{lemma}[An elementary triangular covering radius]
\label{rr:lem:QL}
For every real $L>0$,
\[
 Q_L=\conv\{(-1,0),(1,0),(0,L)\}
 \quad\Longrightarrow\quad
 \mu(Q_L,\Z^2)=\frac12+\frac1L.
\]
\end{lemma}
\begin{proof}
At height $y=f\ge0$, the horizontal section of $\rho Q_L$ is
$|x|\le\rho-f/L$, when the right side is nonnegative.
Given any point modulo $\Z^2$, choose its vertical representative
$f\in[0,1)$ and a horizontal representative with $|x|\le1/2$.
For $\rho=1/2+1/L$, this point belongs to $\rho Q_L$, proving
the upper bound.

Conversely, if $0\le\rho<1/2+1/L$, choose $f\in[0,1)$ with
$\rho<1/2+f/L$.  Consider $(1/2,f)$.  Any nonnegative vertical
residual after subtracting a lattice vector is $f+j$ for an
integer $j\ge0$.  Its horizontal section has half-length at most
$\rho-f/L<1/2$, whereas the horizontal residual of $1/2$ has
absolute value at least $1/2$.  No lattice translate covers this
point, proving the lower bound.
\end{proof}

\begin{theorem}[All real parameters and all ambient dimensions]
\label{rr:thm:alln}
Let $n\ge3$ be an integer and let $r\ge1$ be real.  With two
weights equal to one and the remaining $n-1$ weights equal to $r$,
\begin{equation}\label{rr:eq:alln}
 \boxed{\displaystyle
 \mu_2\bigl(S_n(1,1,r,\ldots,r),\Z^n\bigr)
 =\frac12+\frac1{r+1}.}
\end{equation}
More generally, for any real $0<a\le b$,
\begin{equation}\label{rr:eq:ab}
 \boxed{\displaystyle
 \mu_2\bigl(S_n(a,a,b,\ldots,b),\Z^n\bigr)
 =\frac1{2a}+\frac1{a+b}.}
\end{equation}
\end{theorem}
\begin{proof}
Let $B:\R^n\to\R^2$ be any rank-two quotient with rational
kernel, and put $\Lambda_B=B\Z^n$.  The unweighted columns
\[
 u_0=B(-\one_n),\qquad u_j=Be_j\quad(1\le j\le n)
\]
span the plane and sum to zero.  Consequently zero lies in the
interior of their convex hull.  The image of
$S_n(1,1,r,\ldots,r)$ keeps $u_0,u_1$ fixed and scales all other
columns by $r$.  Theorem~\ref{rr:thm:stretch} therefore gives
\[
 \mu\bigl(BS_n(1,1,r,\ldots,r),\Lambda_B\bigr)\le H_r.
\]
Taking the supremum over all rational rank-two quotients gives
the covering-minimum upper bound, by the standard quotient
characterization \cite[Section 1]{CSS}.

For the reverse inequality, take the integer map
\[
 A(x_1,\ldots,x_n)=(x_2+x_3,x_1).
\]
It satisfies $A\Z^n=\Z^2$, and its image on the simplex is
\[
 P_r=\conv\{(-2,-1),(0,1),(r,0)\}.
\]
The extra vertices for $n>3$ map to zero, which is already
interior to this triangle.  The affine unimodular map
$(x,y)\mapsto(y,x-y+1)$ sends $P_r$ to $Q_{r+1}$.
Covering radius is invariant under translations of the body
and lattice-linear isomorphisms.  Lemma~\ref{rr:lem:QL} gives
$\mu(P_r,\Z^2)=H_r$, proving the matching lower bound.

Finally,
$S_n(a,a,b,\ldots,b)=aS_n(1,1,b/a,\ldots,b/a)$, and covering
minima scale inversely under positive scalar dilation.
\end{proof}

\begin{corollary}[The original weighted conjecture fails throughout the family]
\label{rr:cor:conjecture}
For every $n\ge3$ and every real $r>1$, the value proposed by
Codenotti--Santos--Schymura \cite[Conjecture 5.3]{CSS} for the
second covering minimum of $S_n(1,1,r,\ldots,r)$ is too small.
The exact discrepancy is
\begin{equation}\label{rr:eq:gap}
 \left(\frac12+\frac1{r+1}\right)-\frac{r+2}{2r+1}
 =\frac{r-1}{2(r+1)(2r+1)}>0.
\end{equation}
\end{corollary}
\begin{proof}
The conjecture uses only the first three ordered weights for
$i=2$, and therefore predicts
$(1+2/r)/(2+1/r)=(r+2)/(2r+1)$.  Subtract this value from
\eqref{rr:eq:alln}.
\end{proof}

\begin{corollary}[The complete vector in dimension three]
\label{rr:cor:vector}
For every real $r\ge1$,
\[
 \bigl(\mu_1,\mu_2,\mu_3\bigr)\bigl(S_3(1,1,r,r),\Z^3\bigr)
 =\left(\frac12,\frac{r+3}{2(r+1)},
              \frac{r^2+4r+1}{2r(r+1)}\right).
\]
\end{corollary}
\begin{proof}
The middle value is Theorem~\ref{rr:thm:alln}.  The first and
last are the established endpoint formulas
\cite[Theorem 1.4 and the discussion after Conjecture 5.3]{CSS};
for the last substitute reciprocal weights $(1,1,1/r,1/r)$.
\end{proof}

\begin{remark}[What has been extended]
The parameter $r$ is any real number at least one; no arithmetic
restriction or interpolation from integer cases is used.  The
upper bound controls every rational rank-two quotient.  This
determines $\mu_2$ for the stated family in all dimensions, but
does not assert a formula for arbitrary independently varying
weights or for the intermediate minima $\mu_i$, $i\ge3$, when
$n>3$.  The proof transfers the elementary weighted-flow method
of VMRC [200] \cite{VMRC200}, together with established planar
results, into a new exact covering-minimum calculation.
\end{remark}


\section{An exact finite partition formula for weighted simplices}
\label{part:section}

The jump bound alone does not normally determine a covering radius.
For the relation $(k,1,\ldots,1)$, however, its entire deficit has a
finite combinatorial description.  The weights in this section may
be arbitrary positive real numbers.

Fix integers $d,k\ge1$, positive $a_1,\ldots,a_d,b$, and put
$p=(k,1,\ldots,1)\in\Z^d$.  Define
\begin{equation}\label{part:simplex}
 S(k;a,b)=\conv\left\{-\frac{p}{a_1},\frac{e_1}{b},
              \frac{e_2}{a_2},\ldots,\frac{e_d}{a_d}\right\}.
\end{equation}
Here $a=(a_1,\ldots,a_d)$, and for $d=1$ the displayed list consists
of its first two vertices.  The unscaled directions generate
$\Z^d$.  Listing $e_1$ first, followed by $-p,e_2,\ldots,e_d$,
their primitive positive relation is $(k,1,\ldots,1)$.
Lemma~\ref{tool:gauge} therefore gives
\begin{equation}\label{part:gauge}
 \mu(S(k;a,b),\Z^d)
 =\max_t\min_{s\in[0,1)}
 \left(b\fr{t_0+ks}+\sum_{i=1}^d a_i\fr{t_i+s}\right).
\end{equation}
Write
\begin{equation}\label{part:constants}
 A=\sum_{i=1}^d a_i,\qquad C=kb+A,\qquad
 B=\frac{b+A}{2}-\frac{kb^2+\sum_i a_i^2}{2C}.
\end{equation}
Thus $C$ is the slope of the flow and $B$ is its bound from
Lemma~\ref{tool:jump}.  For $d\ge2$ all the displayed integer
directions are primitive, so $B$ also agrees with the discrete
surface-area expression identified in \cite{VMRC200}.  Our exact
formula remains valid for $d=1$, without requiring that interpretation
of the nonprimitive direction $-k$.

\subsection{The ordered-partition cost}

Let $\mathcal I=(I_0,\ldots,I_{k-1})$ be an ordered partition of
$\{1,\ldots,d\}$; empty parts are allowed.  Set
\begin{align}
 A_h&=\sum_{i\in I_h}a_i,\nonumber\\
 D_0&=0,\qquad
 D_h=\frac{hA}{k}-\sum_{r=0}^{h-1}A_r\quad(1\le h\le k),
       \label{part:levels}\\
 Y_h&=D_h-\min_{0\le r<k}D_r\quad(0\le h<k).\nonumber
\end{align}
In particular $D_k=D_0=0$; extend $Y$ periodically by $Y_k=Y_0$.
In each part, arrange the weights in nonincreasing order,
\[
 a_{h,1}\ge\cdots\ge a_{h,\ell_h},\qquad
 R_{h,j}=\sum_{r=1}^j a_{h,r}.
\]
With $x_+=\max\{x,0\}$, define
\begin{equation}\label{part:cost}
 \mathcal T(\mathcal I)
 =b\sum_{h=0}^{k-1}Y_h
  +\sum_{h=0}^{k-1}\sum_{j=1}^{\ell_h}
        a_{h,j}(Y_h-R_{h,j})_+.
\end{equation}

\begin{theorem}[Exact ordered-partition formula]\label{part:exact}
For the simplex \eqref{part:simplex},
\begin{equation}\label{part:exacteq}
 \boxed{\displaystyle
 \mu(S(k;a,b),\Z^d)
 =B-\frac1C\min_{\mathcal I}\mathcal T(\mathcal I).}
\end{equation}
The minimum ranges over finitely many ordered partitions.  For
rational weights its value can be evaluated using exact rational
arithmetic.
\end{theorem}

\begin{proof}
Call the $k$ crossings of the coordinate with weight $b$
\emph{distinguished jumps}; every other coordinate has a single
jump in a period.  Rotating the variable $s$, place the distinguished
jumps at $h/k$.

We first specify how coincidences are treated.  Resolve simultaneous
crossings in a sequential order, allowing zero time between them.
Every formal intermediate post-jump height is at least the actual
final post-jump height.  Consequently the minimum is unchanged by
including these intermediate states.  If $M=\min F$ and $H_e$ is
the formal post-jump height following an individual jump of size
$w_e$, the individual jump-square identity gives
\begin{equation}\label{part:deficit}
 M=B-\frac1C\sum_e w_e(H_e-M),\qquad H_e-M\ge0.
\end{equation}
This is also \eqref{tool:deficit}.  For example, at coincident jumps
$u$ followed by $v$, the intermediate height is the final height
plus $v$.  Its contribution adds $uv$ to the weighted height sum,
exactly accounting for the cross term between $u^2+v^2$ and
$(u+v)^2$.  Thus individual squared jump sizes must be paired with
these formal heights, rather than with repeated copies of the same
right-continuous value.

\emph{Lower bound on the loss.}
At a distinguished crossing, put that jump first in the tie order.
Assign the remaining singleton jumps to the following half-open
interval $[h/k,(h+1)/k)$.  Their index sets form an ordered
partition $\mathcal I$.  Let $L_h$ denote the formal height
immediately after the distinguished jump at $h/k$.  Across the
interval,
\[
 L_{h+1}-L_h=\frac Ck-A_h-b=\frac Ak-A_h.
\]
Hence $L_h-M=D_h+u$ for a constant $u$.  Since every such excess
is nonnegative, $u\ge-\min_rD_r$, and therefore
$L_h-M\ge Y_h$.

Consider the actual order of singleton jumps in this part.  After
a prefix of total weight $R$, its post-jump excess is at least
$(Y_h-R)_+$: the elapsed time can only increase the height, and
every post-jump excess is nonnegative.  Substituting these estimates
in \eqref{part:deficit} bounds the loss below by the expression
\eqref{part:cost} with the actual within-part orders.

Nonincreasing order minimizes that expression.  Indeed, for adjacent
weights $a\ge c$ and the level $X$ remaining after their common
prefix, the required exchange inequality is
\begin{equation}\label{part:exchange}
 a(X-a)_++c(X-a-c)_+
 \le c(X-c)_++a(X-a-c)_+.
\end{equation}
The right side minus the left side equals, respectively,
\[
 \begin{cases}
 0,&X\le c,\\
 c(X-c),&c\le X\le a,\\
 (a-c)(a+c-X),&a\le X\le a+c,\\
 0,&X\ge a+c.
 \end{cases}
\]
All later terms have the same prefix total after this exchange.
Repeated exchanges therefore prove the assertion.  For every phase
choice we obtain $M\le B-\mathcal T(\mathcal I)/C$, with the
descending orders prescribed in \eqref{part:cost}.  This proves
the upper bound in \eqref{part:exacteq}.

\emph{Attaining the cost of every ordered partition.}
Fix $\mathcal I$ and its descending within-part orders.  Construct a
periodic profile of slope $C$, recording excess above a baseline
that will be chosen below.  At the distinguished jump beginning
part $h$, prescribe excess $Y_h$.  Define
\[
 Z_{h,0}=Y_h,\qquad Z_{h,j}=(Y_h-R_{h,j})_+.
\]
Before singleton jump $j$, wait for
\begin{equation}\label{part:gaps}
 g_{h,j}=\frac{a_{h,j}+Z_{h,j}-Z_{h,j-1}}C\ge0.
\end{equation}
This gives precisely the prescribed post-jump excess $Z_{h,j}$.
The numerator is nonnegative because the positive-part function
decreases by at most $a_{h,j}$ when its argument decreases by that
amount.  After the last singleton, wait for
\begin{equation}\label{part:lastgap}
 g_{h,*}=\frac{b+Y_{h+1}-Z_{h,\ell_h}}C>0
\end{equation}
before the next distinguished jump.  For an empty part use
$Z_{h,\ell_h}=Y_h$.  If $Y_h\ge A_h$, the numerator in
\eqref{part:lastgap} is $b+A/k$; otherwise it is $b+Y_{h+1}$.
This proves its strict positivity.  The total interval length is
\[
 \sum_jg_{h,j}+g_{h,*}
 =\frac{A_h+b+Y_{h+1}-Y_h}{C}
 =\frac{b+A/k}{C}=\frac1k.
\]
Thus the distinguished jumps occur at the required equally spaced
times, and every singleton occurs once in the period.

All prescribed excesses are nonnegative.  Moreover, some $Y_h=0$.
If that part is nonempty, its first waiting time is
$a_{h,1}/C>0$; if it is empty, its final waiting time is positive.
The level zero is therefore an actual right-continuous post-jump
value, followed by a positive interval without a crossing.  It is
not merely an intermediate state hidden inside a simultaneous
event.  The minimum excess is exactly zero, and the weighted sum
of all formal post-jump excesses is exactly $\mathcal T(\mathcal I)$.

The resulting profile has the same slope and jump pattern as a flow
in \eqref{part:gauge}, so it differs from such a flow only by a
constant.  Its mean is determined by the periodic square identity:
if its baseline is $M$, then
\[
 \int_0^1F(s)\,ds
 =M+\frac{\mathcal T(\mathcal I)}C
       +\frac{kb^2+\sum_i a_i^2}{2C}.
\]
Choose $M=B-\mathcal T(\mathcal I)/C$.  This mean is then
$(b+A)/2$, the mean of every fractional-part flow under consideration.
The additive constant is zero.  The constructed flow attains the
claimed minimum for this partition; minimizing the cost completes
the proof.
\end{proof}

\begin{remark}[Computational scope]\label{part:scope}
There are $k^d$ assignments before removing symmetries, so the theorem
does not assert a polynomial-time algorithm.  General finite
algorithms for polytope covering radii and earlier hardness results
are already part of the literature; see \cite{CGNS}.  Here the
additional conclusion is a specific exact partition cost, with no
continuous phase optimization remaining.  When $k=1$ the cost
vanishes, recovering the all-one-relation weighted formula of
\cite[Theorem 1.4]{CSS}.
\end{remark}

\subsection{The exact two-part discrepancy}

For positive $a_1,\ldots,a_d$, define
\begin{equation}\label{part:discrepancy}
 \Delta(a)=\min_{I\subseteq\{1,\ldots,d\}}
       \left|A-2\sum_{i\in I}a_i\right|.
\end{equation}

\begin{theorem}[Two distinguished jumps]\label{part:two}
For $k=2$ and arbitrary positive real weights,
\begin{align}
 \mu(S(2;a,b),\Z^d)
 &=B-\frac{b\Delta(a)}{2(2b+A)}\label{part:twoeq}\\
 &=\frac{b(3A-\Delta(a))+2\sum_{i<j}a_i a_j}
           {2(2b+A)}.\nonumber
\end{align}
\end{theorem}

\begin{proof}
For two parts of weights $A_0,A_1$, the distinguished post-jump
heights differ by $(A_1-A_0)/2$.  Their nonnegative excesses above
the global minimum therefore sum to at least $|A_0-A_1|/2$.
Equation~\eqref{part:deficit} gives a loss of at least
$b|A_0-A_1|/(2C)$, and hence at least $b\Delta(a)/(2C)$.

Choose a minimum-discrepancy partition and label the heavier part
$I_0$.  Set $\delta=\Delta(a)/2$, so that
$A_0=A/2+\delta$ and $A_1=A/2-\delta$.  If $\delta>0$, the
heavier part contains an element of weight at least $\delta$:
otherwise moving any one of its positive weights $a<\delta$ to
the other part would change the discrepancy from $2\delta$ to
$2\delta-2a$.  If $\delta=0$, either part contains a positive
weight because $A>0$.

Prescribe distinguished post-jump excesses $\delta$ and $0$ at
times $0$ and $1/2$.  In the heavy part, place a weight
$a\ge\delta$ first after time $(a-\delta)/C$, bringing the excess
to zero.  Place each remaining singleton after a gap equal to its
weight divided by $C$, keeping every singleton post-jump excess
zero.  The final gap before the next distinguished jump is $b/C$.
The total duration is
\[
 \frac{A_0-\delta+b}{C}=\frac12.
\]
In the light part start at zero, use singleton gaps $a_i/C$, and
end with a gap $(b+\delta)/C$.  Its total duration is also $1/2$,
and the final distinguished jump returns to excess $\delta$.
An empty light part is allowed.  Zero initial waiting time is
handled by the formal coincidence convention in the preceding proof.

The weighted post-jump excess sum is exactly $b\delta$.  Choosing
the baseline $B-b\delta/C$ and matching means as in
Theorem~\ref{part:exact} makes this profile a fractional-part flow.
It attains the upper bound.  Expanding $B$ gives the second formula.
\end{proof}

For two singleton weights $a,c$, the discrepancy is $|a-c|$, and
\begin{equation}\label{part:twosingletons}
 \mu=\frac{b(3(a+c)-|a-c|)+2ac}{2(2b+a+c)}.
\end{equation}
If $a\ge c$, this is $(ba+2bc+ac)/(2b+a+c)$.  In particular,
the exact value is a rational expression on each of the two order
regions, and equality in the flow bound holds precisely when $a=c$.

\begin{figure}[ht]
\centering
\begin{tikzpicture}[x=9.8cm,y=1.15cm]
 \draw[->,gray] (-.035,0)--(1.07,0) node[right]{$s$};
 \draw[->,gray] (0,0)--(0,3.65) node[above]{$F(s)$};
 \draw[dotted,accent] (0,1.2)--(1.03,1.2);
 \draw[very thick,accent] (0,1.7)--(.3,3.2);
 \draw[very thick,accent] (.3,1.2)--(.5,2.2);
 \draw[very thick,accent] (.5,1.2)--(.7,2.2);
 \draw[very thick,accent] (.7,1.2)--(1,2.7);
 \foreach \x/\y/\z in {.3/3.2/1.2,.5/2.2/1.2,.7/2.2/1.2,1/2.7/1.7}{
  \draw[->,gray] (\x,\y)--(\x,\z);
  \draw[accent,fill=white] (\x,\y) circle(1.8pt);
  \fill[accent] (\x,\z) circle(1.8pt);}
 \fill[accent] (0,1.7) circle(1.8pt);
 \foreach \x/\lab in {0/0,.3/{3/10},.5/{1/2},.7/{7/10},1/1}{
  \node[below,font=\small] at (\x,0){$\lab$};}
 \node[left,font=\small] at (0,1.2){$6/5$};
 \node[left,font=\small] at (0,1.7){$17/10$};
 \node[left,font=\small] at (0,3.2){$16/5$};
\end{tikzpicture}
\caption{An attained partition deficit. For $k=2$, singleton weights
$(2,1)$ and repeated weight $b=1$, the flow is
$\{2s\}+2\{s-3/10\}+\{s-7/10\}$.
Its minimum is $6/5$, while $B=13/10$ and the minimum partition
discrepancy is $1$. The exact loss is $1/10$.}
\label{fig:partitiondeficit}
\end{figure}

\subsection{Equal singleton weights and a greatest-common-divisor formula}

\begin{theorem}[Equal singleton weights]\label{part:gcd}
If $a_1=\cdots=a_d=a>0$ and $g=\gcd(d,k)$, then
\begin{equation}\label{part:gcdeq}
 \boxed{\displaystyle
 \mu(S(k;a\one,b),\Z^d)
 =\frac{(d-1)a}{2}+\frac{ab(d+g)}{2(kb+da)}.}
\end{equation}
\end{theorem}

\begin{proof}
For any ordered partition write $A_h=an_h$, with nonnegative
integers $n_h$ summing to $d$.  If $D_m$ is a minimum, then
\[
 \frac{Y_h}{a}
 =\frac{(h-m)d}{k}-\text{an integer}.
\]
The fractional parts of these $k$ quantities run through
$0,g/k,\ldots,(k-g)/k$, each $g$ times, and sum to $(k-g)/2$.
Since $Y_h\ge0$, this gives
\[
 \sum_hY_h\ge\frac{a(k-g)}2,
 \qquad \mathcal T(\mathcal I)\ge\frac{ab(k-g)}2.
\]
For the reverse inequality choose the part sizes
\[
 n_h=\left\lfloor\frac{(h+1)d}{k}\right\rfloor
       -\left\lfloor\frac{hd}{k}\right\rfloor.
\]
Then $D_h=a\fr{hd/k}$, so $Y_h=D_h\in[0,a)$.  Every nonempty
prefix has weight at least $a$, and therefore every singleton term
in \eqref{part:cost} vanishes.  The cost is exactly $ab(k-g)/2$.
Substitution in \eqref{part:exacteq} and simplification prove
\eqref{part:gcdeq}.
\end{proof}

At $a=b=1$ this gives the unweighted formula
\begin{equation}\label{part:unweighted}
 \mu\bigl(\conv\{-(k,1,\ldots,1),e_1,\ldots,e_d\}\bigr)
 =\frac{d-1}{2}+\frac{d+\gcd(d,k)}{2(d+k)}.
\end{equation}
The instance $d=3,k=2$ is $7/5$, already computed in
\cite[Corollary B.4]{CSS}; the new statement here is the general
weighted formula, rather than that isolated value.

\subsection{Equality and the complexity of attaining the bound}

\begin{corollary}[Balanced-partition criterion]\label{part:balanced}
For arbitrary positive real weights,
\begin{equation}\label{part:balancedeq}
 \mu(S(k;a,b),\Z^d)=B
 \quad\Longleftrightarrow\quad
 \{1,\ldots,d\}\text{ has a partition into $k$ parts of weight }A/k.
\end{equation}
The existence of equality is independent of the positive weight $b$.
For equal singleton weights it holds if and only if $k$ divides $d$.
\end{corollary}

\begin{proof}
Every term in \eqref{part:cost} is nonnegative and $b>0$.
Thus the cost is zero if and only if all $Y_h$ vanish.  By
\eqref{part:levels}, this is equivalent to $A_h=A/k$ for every $h$.
The final assertion follows because equal weights require every
part to contain exactly $d/k$ elements.
\end{proof}

\begin{corollary}[A restricted equality problem is strongly NP-complete]
\label{part:complexity}
With dimension allowed to vary, deciding whether the rational
simplex \eqref{part:simplex} attains its explicit rational bound
$B$ is strongly NP-complete, already on the following restricted
inputs: $b=1$, $d=3k$, and positive integers $a_i,B_0$ satisfying
\[
 \frac{B_0}{4}<a_i<\frac{B_0}{2},\qquad
 \sum_i a_i=kB_0.
\]
In particular the hardness holds with all integer parameters and
all denominators of the simplex vertices bounded polynomially in
the input dimension.
\end{corollary}

\begin{proof}
Use the standard strongly NP-complete problem
\textsc{3-Partition} with precisely these promises \cite{GJ}.
By Corollary~\ref{part:balanced}, equality holds exactly when the
$3k$ weights can be divided into $k$ groups of sum $B_0$.
The strict bounds force each such group to have exactly three
elements.  This is exactly the original instance, and the mapping
to \eqref{part:simplex} and \eqref{part:constants} has polynomial
encoding length.  The strong restriction of \textsc{3-Partition}
keeps $B_0$, and hence all $a_i$ and all vertex denominators,
polynomially bounded.  Conversely a partition is a certificate
verifiable by exact rational arithmetic in polynomial time.
Thus the restricted equality problem belongs to NP and is strongly
NP-hard.
\end{proof}

This conclusion concerns equality in a specified bound on a specified
simplex family.  General hardness of covering-radius computation
predates it, as discussed in \cite{CGNS}; no claim of first hardness
for lattice covering is needed.

\section{A finite linear classification of weighted equality}
\label{part:wordssection}

The partition criterion is a special case of a finite description
for any fixed positive primitive relation.  Let
$u_1,\ldots,u_m$ generate a rank-$(m-1)$ lattice $\Lambda$, and let
their column matrix have kernel $\R q$, where
$q=(q_1,\ldots,q_m)\in\Z_{>0}^m$ is primitive.  For
$c\in\R_{>0}^m$, put
\[
 P(c)=\conv\{u_1/c_1,\ldots,u_m/c_m\},\qquad
 C=\sum_jq_jc_j,\qquad
 B(q,c)=\frac12\sum_jc_j-\frac{\sum_jq_jc_j^2}{2C}.
\]

A cyclic word $W$ is admissible in multiplicity if it contains
exactly $q_j$ copies of symbol $j$.  For two consecutive occurrences
of a symbol $j$, count the symbols strictly after the first and
up to and including the second, following the cyclic order.
Write this count vector as $N=(N_1,\ldots,N_m)$.  When $q_j=1$,
the interval is one full traversal of the word.

\begin{theorem}[Cyclic-word equality criterion]\label{part:word}
The equality $\mu(P(c),\Lambda)=B(q,c)$ holds if and only if
there is an admissible cyclic word for which every consecutive
pair of occurrences of every symbol $j$ satisfies
\begin{equation}\label{part:wordeq}
 q_j\sum_{\ell=1}^m N_\ell c_\ell
       =\sum_{\ell=1}^m q_\ell c_\ell.
\end{equation}
For each fixed $q$, the equality weights form a finite union of
intersections of rational linear subspaces with $\R_{>0}^m$.
\end{theorem}

\begin{proof}
By Lemma~\ref{tool:gauge}, the covering radius is the maximum of
the minimum of $\sum_jc_j\fr{t_j+q_js}$, and that maximum is
attained.  Suppose it equals $B(q,c)$.  Lemma~\ref{tool:jump}
then implies that all crossings are distinct and all post-jump
heights coincide.  List their coordinate labels cyclically to
obtain $W$.  The gap preceding an event of label $\ell$ is
$c_\ell/C$, because slope $C$ must exactly compensate its jump.
Consecutive occurrences of coordinate $j$ are separated by
$1/q_j$.  Summing the intervening gaps gives
$\sum_\ell N_\ell c_\ell/C=1/q_j$, which is
\eqref{part:wordeq}.

Conversely, take an admissible word satisfying these equations.
Lay its events on a circle of circumference one, with the gap
preceding an event of label $\ell$ equal to $c_\ell/C$.
The gaps are positive and their sum is one.  Equation
\eqref{part:wordeq} says that successive events with label $j$
are separated by $1/q_j$.  Thus the events of every label are
exactly the crossing times of a coordinate
$\fr{t_j+q_js}$ for an appropriate phase $t_j$.
Between each pair of events the increase $C$ times the gap is
exactly the next jump size.  All post-jump values of the resulting
flow therefore coincide.  Lemma~\ref{tool:jump} gives equality,
and Lemma~\ref{tool:gauge} transfers it to the covering radius.

For a fixed word, \eqref{part:wordeq} is a finite homogeneous
linear system with integer coefficients in $c$.  There are finitely
many words of the prescribed multiplicities, proving the final
statement.
\end{proof}

\begin{corollary}[Generic strictness]\label{part:measurezero}
If the primitive relation $q$ is not $(1,\ldots,1)$, the set of
positive weights attaining $B(q,c)$ has Lebesgue measure zero in
$\R^m$.  If $q=(1,\ldots,1)$, every positive weight vector
attains the bound.
\end{corollary}

\begin{proof}
In the first case choose $j$ with $q_j>1$.  If the linear system
associated with a word imposed no nonzero equation, one of its
consecutive-$j$ equations would have coefficients
$q_jN_\ell-q_\ell=0$ for every $\ell$.  Then $q_j$ would divide
every coordinate of $q$, contradicting primitiveness.  Every word
therefore defines a proper linear subspace.  A finite union of
proper linear subspaces has measure zero.

If all $q_j=1$, each consecutive-symbol interval traverses the
whole word, so $N_\ell=1$ for every $\ell$.  All equations are
identities for every positive $c$; Theorem~\ref{part:word} applies.
\end{proof}

\begin{remark}[Scope of the classification]\label{part:wordscope}
The word criterion also characterizes equality for the scalar
fractional-part optimization itself.  Its geometric equivalence
above uses a simplex, a one-dimensional positive kernel, and
generation of the ambient lattice.  It does not assert the same
necessity for an arbitrary configuration with a higher-dimensional
kernel.  The classification is finite, but its number of words
can be large; it is not an efficiency claim or an automatic
solution of unrelated covering conjectures.
\end{remark}


\section{A four-dimensional equality counterexample}
For completeness, we retain the explicit equality counterexamples
from the first transfer note. They strengthen the dimensional
counterexample in VMRC [200]. Let
\begin{equation}\label{eq:S4}
 S=\conv\{(-2,-1,-1,-1),e_1,e_2,e_3,e_4\}\subset\R^4.
\end{equation}

\begin{theorem}[Smallest possible counterexample dimension]\label{thm:S4}
The simplex $S$ has $0\in\inter S$, normalized volume $6$, and
$\mu(S,\Z^4)=2$. It is not obtained from standard terminal simplices
by translations and direct sums, up to unimodular equivalence.
It therefore disproves the equality clause of Conjecture A in
dimension four. This is the smallest possible dimension.
\end{theorem}
\begin{proof}
Write $v_0=(-2,-1,-1,-1)$, $v_j=e_j$ for $j\ge1$. The positive
primitive relation is $q=(1,2,1,1,1)$, with sum $Q=6$. Its normalized
coefficients give strictly positive barycentric coordinates for zero.
The standard basis vertices generate the whole lattice.

The residue classes
\[
 1\pmod6,\quad0\pmod3,\quad2\pmod6,\quad4\pmod6,\quad5\pmod6
\]
partition the integers. Choose
\[
 t=(-1/6,0,-1/3,-2/3,-5/6).
\]
There is exactly one crossing at each $s=j/6$, and
\begin{equation}\label{eq:S4cert}
 \sum_{j=0}^4\{t_j+q_js\}=2+\{6s\}\qquad(s\in\R).
\end{equation}
Both sides have slope $6$, the same unit jumps, and value $2$ at zero,
which proves the all-real identity including its endpoints.
Lemma~\ref{tool:gauge} gives the global lower bound $2$ at
\[
 x_*=(1/3,-1/6,-1/2,-2/3).
\]
Lemma~\ref{tool:jump}, with all $c_j=1$, gives the matching upper bound
for every point.

Using edges from $v_0$, the normalized volume is
\[
 \det(I_4+(2,1,1,1)^{\mathsf T}\one_4^{\mathsf T})=6.
\]
A standard terminal $j$-simplex has normalized volume $j+1$.
Translations preserve volume and direct sums multiply normalized
volumes. Thus a body in the proposed family of dimension four has
volume in
\[
 \{5,8,9,12,16\},
\]
as obtained from the five integer partitions of four. Volume six is
excluded. These invariants are unchanged by unimodular maps.
Finally, the complete equality classification is already proved
in dimensions at most three \cite[Theorem 1.3]{CSS}.
\end{proof}

For clarity, the volume multiplication used here remains valid when
the origin is on the boundary of a summand. For bodies of dimensions
$a,b$, the fiber formula for their free sum gives
\[
 \vol_{a+b}(K\oplus L)
 =\frac{a!b!}{(a+b)!}\vol_a(K)\vol_b(L),
\]
by integrating the scaled fibers $(1-g_K(x))L$ over $K$.
The sublevel volume of the possibly extended gauge $g_K$ is
$t^a\vol_a(K)$; the beta integral proves the displayed identity.
Multiplication by $(a+b)!$ gives normalized-volume multiplication.

\section{Equality counterexamples with every prescribed \texorpdfstring{$k\ge2$}{k >= 2}}
Conjecture D of \cite{CSS} proposes, for a lattice $d$-polytope with
$k\ge1$ interior lattice points,
\[
 \mu(P)\le\frac{d-1}{2}+\frac1{k+1},
\]
and specifies an equality family. For $k\ge2$ that family is formed
from a segment of length $k+1$ and standard terminal simplices by
translations and direct sums. The following result concerns its
\emph{equality clause}, not the general numerical upper bound.

\begin{theorem}[Five-dimensional counterexamples for all $k\ge2$]
\label{thm:D}
For the simplex $S$ in \eqref{eq:S4}, put
\[
 R_k=S\oplus[-1,k]\subset\R^4\oplus\R,\qquad k\in\Z,\ k\ge2.
\]
Then
\[
 |\inter R_k\cap\Z^5|=k,\qquad
 \mu(R_k,\Z^5)=2+\frac1{k+1},\qquad
 \Vol_{\Z^5}(R_k)=6(k+1).
\]
No $R_k$ belongs to the equality family proposed in Conjecture D.
\end{theorem}
\begin{proof}
The facet inequalities of $S$ are $\ell_j(x)\le1$, where
\begin{align*}
 \ell_0(x)&=x_1+x_2+x_3+x_4,\\
 \ell_1(x)&=-2x_1+x_2+x_3+x_4,\\
 \ell_2(x)&=x_1-5x_2+x_3+x_4,\\
 \ell_3(x)&=x_1+x_2-5x_3+x_4,\\
 \ell_4(x)&=x_1+x_2+x_3-5x_4.
\end{align*}
They satisfy $\ell_0+2\ell_1+\ell_2+\ell_3+\ell_4=0$.
For an interior lattice point, all these integer-valued forms are
nonpositive, so positivity of the relation forces all to be zero;
their common kernel is zero. Thus zero is the only interior lattice
point of $S$. More precisely its gauge
$g_S(z)=\max_j\ell_j(z)$ is an integer at every $z\in\Z^4$ and
is at least one when $z\ne0$.

The gauge of a free sum is the sum of the two gauges. A lattice point
$(z,t)$ can therefore be in $\inter R_k$ only if $z=0$. For $z=0$,
the permitted integers are precisely $-1<t<k$, namely
$0,1,\ldots,k-1$. This proves the count.

For a direct-product lattice, minimizing the sum gauge over lattice
translates separates in the two factors. Taking the two independent
suprema yields additivity of covering radii, also stated in
\cite[Corollary 2.2]{CSS}. Hence
$\mu(R_k)=\mu(S)+\mu([-1,k])=2+1/(k+1)$.
The normalized-volume identity gives $6(k+1)$.

If $R_k$ were in the proposed family, the terminal-simplex factors
would have dimensions with total four and volumes with product six,
after canceling the common segment factor $k+1$. The possible
products in dimension four are $5,8,9,12,16$, none equal to six.
This excludes all permitted translations, direct-sum groupings and
unimodular images.
\end{proof}



\section{The numerical interior-point bound: an exact family and a closure principle}
\label{d:section}

The equality clauses discussed above are distinct from the numerical
part of CSS Conjecture~D \cite{CSS}. For a lattice $d$-polytope with
$k\ge1$ interior lattice points, that part asks whether
\begin{equation}\label{d:conjecture}
 \mu(P,\Lambda)\le\frac{d-1}{2}+\frac1{k+1}.
\end{equation}
We do not settle \eqref{d:conjecture} for arbitrary polytopes.
We determine the radius and the interior-point count for an infinite
family in every dimension, and establish a direct-sum closure rule.

\begin{theorem}[An exact family in every dimension]\label{d:family}
For integers $d,n\ge1$, let
\[
 S_{d,n}=\conv\bigl\{-(1,\ldots,1,n),e_1,\ldots,e_d\bigr\}
 \subset\R^d,
 \qquad g=\gcd(d,n).
\]
Then
\begin{equation}\label{d:exact}
 \bigl|\inter S_{d,n}\cap\Z^d\bigr|
 =k:=\left\lceil\frac nd\right\rceil,
 \qquad
 \mu(S_{d,n},\Z^d)
 =\frac{d-1}{2}+\frac{d+g}{2(d+n)}.
\end{equation}
In particular, \eqref{d:conjecture} holds for every $S_{d,n}$.
When $k\ge2$, equality holds if and only if $d\mid n$.
\end{theorem}
\begin{proof}
The primitive positive vertex relation has $d$ entries equal to
one and one entry equal to $n$, with total $d+n$. The vertices
generate $\Z^d$. The radius formula is the specialization $a=b=1$
of Theorem~\ref{part:gcd}, with one frequency $n$ and $d$ unit
frequencies, all carrying unit weight. It therefore applies also
when $n$ and $d$ are not coprime.

For $i<d$, the $i$th coordinate of an interior lattice point must
be zero, since it is strictly between $-1$ and $1$. Such a point
has the form $t e_d$, where $t\in\Z$. Its barycentric coordinates
in the displayed vertex order are
\[
 \alpha_0=\cdots=\alpha_{d-1}=\frac{1-t}{d+n},
 \qquad \alpha_d=\frac{n+dt}{d+n}.
\]
They are all positive exactly when $-n/d<t<1$. Thus the interior
points are $0,-e_d,\ldots,-(k-1)e_d$, proving the count. This
argument also covers $d=1$, where $S_{1,n}=[-n,1]$.

Write $n=d(k-1)+r$, with $1\le r\le d$. Then
$d+n=dk+r$ and $g=\gcd(d,r)\le r$. Direct subtraction gives
\begin{equation}\label{d:gap}
 \frac{d-1}{2}+\frac1{k+1}-\mu(S_{d,n})
 =\frac{(k-1)(d-g)+2(r-g)}{2(dk+r)(k+1)}\ge0.
\end{equation}
For $k\ge2$, the numerator vanishes exactly when $g=d$ and $r=g$,
equivalently $r=d$, or $d\mid n$.
\end{proof}

\begin{remark}[Established special cases and the present scope]
The case $n=1$ is the standard terminal-simplex radius already
computed in the 2016 preprint \cite[Proposition~4.8]{GMS}. The value
$\mu(S_{3,2})=7/5$ appears in \cite[Corollary~B.4]{CSS}.
Theorem~\ref{d:family} supplies a uniform exact formula and the
interior-count comparison for all $d,n$; these previously known
special cases are not new conjecture resolutions. No claim of
historical priority for the general formula is made here.
\end{remark}

\begin{proposition}[Closure under direct sums]\label{d:closure}
Let $P_i$ be lattice polytopes of positive dimensions $d_i$, with
$0\in\inter P_i$ and $k_i\ge1$ interior lattice points, for $i=1,2$.
If each $P_i$ satisfies \eqref{d:conjecture}, then their direct sum
in the direct-sum lattice also satisfies \eqref{d:conjecture}.
If both $k_1$ and $k_2$ exceed one, the resulting inequality is strict.
\end{proposition}
\begin{proof}
The gauge of $P_1\oplus P_2$ is the sum of the factor gauges.
Every interior lattice pair therefore has both components in the
respective interiors. If $k$ is the number of interior lattice
points of the direct sum, then $k\le k_1k_2$.
Minimizing the sum gauge over the direct-product lattice separates
in the two factors, and taking the independent suprema gives
additivity of covering radii; see also \cite[Corollary~2.2]{CSS}.
Consequently
\[
 \mu(P_1\oplus P_2)
 \le\frac{d_1+d_2-2}{2}+\frac1{k_1+1}+\frac1{k_2+1}.
\]
For all real $a,b\ge1$, the exact identity
\[
 \frac12+\frac1{ab+1}-\frac1{a+1}-\frac1{b+1}
 =\frac{(a-1)(b-1)(ab-1)}{2(a+1)(b+1)(ab+1)}\ge0
\]
holds. Apply it to $a=k_1$, $b=k_2$, and use
$1/(k_1k_2+1)\le1/(k+1)$. This proves the assertion, including
strictness when both counts exceed one.
\end{proof}

\begin{corollary}\label{d:generatedfamilies}
The numerical bound \eqref{d:conjecture} holds for every iterated
direct sum, with interior origins, of polytopes selected from the
following classes: lattice polygons; the simplices $S_{d,n}$;
and lattice polytopes with exactly one interior lattice point.
\end{corollary}
\begin{proof}
For polygons with at least two interior points the assertion is
\cite[Theorem~6.3]{CSS}; the one-point case follows from
Proposition~\ref{tool:nonhollow}. For the second class use
Theorem~\ref{d:family}, and for the third use
Proposition~\ref{tool:nonhollow}. Apply
Proposition~\ref{d:closure} inductively.
\end{proof}

\section{A global qualitative consequence: every maximizer has one interior point}
\label{d:unique-section}

The numerical gap in Conjecture~D is still unresolved in general.
Nevertheless, the equality mechanism of the fractional-part bound
gives a global qualitative restriction, including for ambient
lattices strictly finer than the vertex-generated lattice.

\begin{lemma}[The simplex case in the full ambient lattice]\label{d:simplex-unique}
Let $S$ be a lattice $h$-simplex in a lattice $\Lambda$, with
$0\in\inter S$ and $\mu(S,\Lambda)=h/2$. Then its facet equations
can be written as $\ell_j(x)\le1$, with each $\ell_j$ integral
on $\Lambda$. In particular, $0$ is its only interior lattice point.
\end{lemma}
\begin{proof}
Write $S=\conv(v_0,\ldots,v_h)$, and let
$\Lambda_0=\sum_j\Z v_j\subseteq\Lambda$. If $q$ is the primitive
positive relation among the vertices and $Q=\sum_jq_j$, then
Lemma~\ref{tool:gauge} with unit weights and
Lemma~\ref{tool:jump} give $\mu(S,\Lambda_0)\le h/2$.
Lattice monotonicity gives the reverse inequality, hence
$\mu(S,\Lambda_0)=h/2$.

Consider any deepest point $x$ for $\Lambda_0$, and represent it
as $x=\sum_jt_jv_j$. Its fractional-part function
$F(s)=\sum_j\{t_j+q_js\}$ has minimum $h/2$. Equality must hold
throughout the jump-square estimate. Since all averaging weights
are positive, every post-jump value is $h/2$; equality in the
square-jump estimate excludes simultaneous coordinate crossings.
Thus every jump has size one, and the constant slope $Q$ forces
all successive event gaps to equal $1/Q$.

Shift $s$ so that the event set is $\{a/Q:0\le a<Q\}$.
Coordinate $j$ has events spaced by $1/q_j$, so
$m_j=Q/q_j$ is an integer. Its events form one residue class
$a_j\pmod {m_j}$, with $0\le a_j<m_j$. These residue classes
partition the event grid. Exactly one $a_j$ is zero, and modulo
$\Lambda_0$ the point $x$ has the coefficient representation
$t_j=-a_j/m_j$. Evaluating $F$ at zero gives
\[
 \frac h2=h-\sum_j\frac{a_jq_j}{Q},
 \qquad\text{so}\qquad
 \sum_j\frac{a_jq_j}{Q}=\frac h2.
\]

Define linear facet forms by
\[
 \ell_j(v_i)=1\quad(i\ne j),\qquad
 \ell_j(v_j)=1-Q/q_j.
\]
These values respect the unique linear relation among the vertices,
so the forms are well-defined. They are integral on $\Lambda_0$,
and $S=\{x:\ell_j(x)\le1\text{ for all }j\}$.
The coefficient representation above now proves, for \emph{every}
deepest point of the coarse lattice,
\begin{equation}\label{d:deep-congruence}
 \ell_j(x)\equiv-\frac h2+a_j
 \equiv-\frac h2\pmod\Z.
\end{equation}

Choose a deepest point $x$ for the full lattice $\Lambda$.
Such a point exists because gauge distance is continuous and
lattice-periodic, and hence attains its maximum on the compact
lattice torus. For every $z\in\Lambda$, the distance of $x-z$
to $\Lambda_0$ is at least its distance to $\Lambda$, namely $h/2$,
and at most $\mu(S,\Lambda_0)=h/2$. Thus both $x$ and $x-z$ are
deepest points for $\Lambda_0$. Equation~\eqref{d:deep-congruence}
implies $\ell_j(z)=\ell_j(x)-\ell_j(x-z)\in\Z$.
This proves integrality on the full ambient lattice, without
assuming that the vertices generate that lattice.

Finally $\sum_jq_j\ell_j=0$. If $z\in\inter S\cap\Lambda$,
each integer $\ell_j(z)$ is less than one, hence nonpositive.
The positive relation forces every $\ell_j(z)$ to vanish.
The common kernel of the facet forms is zero: in coefficients
$z=\sum_i t_iv_i$, these equations give
$t_j=(q_j/Q)\sum_i t_i$, so the coefficient vector is a multiple
of $q$ and $z=0$.
\end{proof}

\begin{theorem}[Unique interior point at the maximal radius]\label{d:global-unique}
Let $P$ be a non-hollow lattice $d$-polytope in a full-rank lattice
$\Lambda$. If $\mu(P,\Lambda)=d/2$, then $P$ has exactly one
interior lattice point. Consequently, whenever
$|\inter P\cap\Lambda|\ge2$, one has $\mu(P,\Lambda)<d/2$.
\end{theorem}
\begin{proof}
We use Proposition~\ref{tool:nonhollow} and induct on $d$.
Translate an interior lattice point to zero. Choose an
inclusion-minimal set of vertices of $P$ whose convex hull contains
zero. Its convex hull $S$ is a simplex with $0\in\operatorname{relint}S$.
Put $r=\dim S\ge1$ and $L=\operatorname{span}S$.

If $r=d$, then
$d/2=\mu(P,\Lambda)\le\mu(S,\Lambda)\le d/2$.
Every simplex is tight for every lattice
\cite[Lemma~2.7]{CSS}: any strictly larger convex body in the same
space has strictly smaller covering radius. Therefore $P=S$,
and Lemma~\ref{d:simplex-unique} proves the assertion.

Suppose $r<d$. The subspace $L$ is rational; give it the lattice
$\Lambda\cap L$, and give the quotient the projected lattice
$\pi\Lambda$. Lemma~\ref{tool:section} and the non-hollow bounds give
\[
 \frac d2=\mu(P,\Lambda)
 \le\mu(P\cap L,\Lambda\cap L)+\mu(\pi P,\pi\Lambda)
 \le\frac r2+\frac{d-r}{2}=\frac d2.
\]
Here the section bound follows from $S\subseteq P\cap L$ and
$\mu(S,\Lambda\cap L)\le r/2$. Thus every inequality is an equality.
In particular, tightness of the full-dimensional simplex $S$
\emph{within $L$} forces $P\cap L=S$. This application permits
$P\cap L$ to have nonlattice vertices: the tightness result applies
to all convex super-bodies, not only lattice polytopes.

The quotient $\pi P$ is a non-hollow lattice $(d-r)$-polytope,
with radius $(d-r)/2$. By induction its unique interior lattice
point is zero. A linear surjection takes interior points to
interior points, so every interior lattice point of $P$ projects
to zero. It therefore lies in $L$, and hence in the relative
interior of $P\cap L=S$. Lemma~\ref{d:simplex-unique}, applied in
$L$, says that this point is zero.

The last assertion follows by contraposition together with
$\mu(P,\Lambda)\le d/2$.
\end{proof}

\begin{remark}[The remaining quantitative gap]\label{d:remaining-gap}
Theorem~\ref{d:global-unique} is a strict inequality with no
specified uniform improvement. Conjecture~D asks for the particular
gap $1/2-1/(k+1)$ below $d/2$, which is at least $1/6$ for $k\ge2$.
The theorem does not supply that gap.
One cannot insert it by assuming that a line with $k$ interior
lattice points has a section of lattice length at least $k+1$.
For example, $S_{3,4}$ has exactly two interior points, $0,-e_3$,
but its section on their line has endpoints $-4e_3/3,e_3$ and
length $7/3<3$. The basic section-and-projection bound gives
$1+3/7=10/7$, larger than the conjectured $4/3$, whereas the exact
formula gives $\mu(S_{3,4})=9/7$. Coupling the two parts, rather
than rounding section endpoints to fictitious lattice points,
is essential in this example.
\end{remark}


\section{Dependence, verification, and remaining questions}
\label{sec:scope}
The methods in this manuscript distinguish three mathematical tasks:
finding an obstruction in one quotient, controlling every quotient,
and computing the exact deficit in a phase-uniform inequality.
The planar stretching theorem performs the second task for two fixed
columns in arbitrary planar configurations. It yields an exact
second covering minimum in every ambient dimension. The ordered
partition formula performs the third task for every positive real
weight vector when the relation is $(k,1,\ldots,1)$.

The elementary jump balance and gauge transfer come from VMRC [200].
Their required forms, including the non-hollow bound, have been
proved again in Section~\ref{sec:tools}. Established inputs from
CSS are explicitly cited: the quotient characterization, the planar
interior-point theorem, the weighted terminal-simplex radius formula,
and tightness of simplices. The independent diagonal-quotient
calculation in the appendix extends a documented special-case
method from the 2016 preprint \cite{GMS}.

The original weighted conjecture is a published statement in
\cite{CSS} and is repeated in the January 2026 preprint
\cite{Krivokuca}. The exact examples above contradict its stated
formula. Historical priority for the full transfer results is not
certified by this manuscript. Searches of the original and recent
primary sources and of the displayed formulas did not identify
an identical general statement as of 20 November 2026; this is a
bounded search report. General algorithms and hardness for covering
radii are prior results \cite{CGNS}. The complexity corollary here
concerns sharpness of a specified rational bound on a restricted
simplex family.

The main new derivations were checked independently. In addition,
a direct event-time linear program, separate from the partition-cost
recurrence, agreed with the formula in eight small cases spanning
870 event-order cells. The exact arithmetic certificate below checks
129 constructed profiles, including 79 with simultaneous crossings,
and the weighted gcd and parameter-gap specializations. These finite
checks supplement the proofs; they are not substitutes for the
all-parameter arguments.

The scope left open is precise. We do not give all covering minima
for arbitrary independently varying weights. For the two-level
family in dimension above three, the exact theorem concerns the
second minimum. Conjecture D's numerical bound for arbitrary
non-hollow lattice polytopes is not proved here: the global strict
inequality $\mu<d/2$ for two or more interior points does not supply
its prescribed quantitative gap. The explicit equality
counterexamples, however, do disprove the original classification
clauses in their full stated domains.

\appendix


\section{An independent diagonal-quotient calculation}
\label{d:graph-appendix}

The next discrete calculation independently recovers the unweighted
case of Theorem~\ref{part:gcd}. It extends the diagonal-quotient
description used for the terminal simplex in the 2016 preprint
\cite[Proposition~4.8 and its proof]{GMS} to arbitrary modulus.

\begin{proposition}\label{d:graph-diameter}
For positive integers $d,Q$, let $G_{d,Q}$ be the Cayley digraph of
$(\Z/Q\Z)^{d-1}$ with generators
$e_1,\ldots,e_{d-1},-\one_{d-1}$. For $d=1$ take the one-vertex
graph. Then
\[
 \operatorname{diam}G_{d,Q}
 =\frac{(d-1)Q-d+\gcd(d,Q)}2.
\]
\end{proposition}
\begin{proof}
Write $[u]_Q$ for the least nonnegative residue, put $g=\gcd(d,Q)$,
and let $a=(a_1,\ldots,a_{d-1})$ be a vertex. Set $a_0=0$.
Commutativity allows a shortest walk to use $s$ negative-diagonal
steps followed by coordinate steps. Reducing all counts modulo $Q$
gives
\[
 \operatorname{dist}(0,a)
 =\min_{0\le s<Q}G_a(s),\qquad
 G_a(s)=\sum_{j=0}^{d-1}[a_j+s]_Q.
\]
The mean of $G_a$ is $d(Q-1)/2$, while its residues modulo $Q$
are $A+ds$, with $A=\sum_j a_j$. They traverse one residue class
modulo $g$, with each of its $Q/g$ representatives appearing
exactly $g$ times. If $M=\min_s G_a(s)$, the least possible
representatives at or above $M$ are
$M,M+g,\ldots,M+Q-g$. Therefore
\[
 \frac{d(Q-1)}2\ge M+\frac{Q-g}{2},
 \qquad M\le\frac{(d-1)Q-d+g}{2}.
\]

For equality, choose $a_j=\lfloor jQ/d\rfloor$ for $0\le j<d$.
For each $0\le s<Q$, exactly $\lfloor ds/Q\rfloor$ of these
summands wrap past $Q$, so
\[
 G_a(s)=A+ds-Q\lfloor ds/Q\rfloor=A+[ds]_Q.
\]
Pairing the floors with indices $j$ and $d-j$ gives
\[
 A=\sum_{j=0}^{d-1}\left\lfloor\frac{jQ}{d}\right\rfloor
 =\frac{(d-1)(Q-1)+g-1}{2}.
\]
Indeed there are $g-1$ nonzero indices with $d\mid jQ$; their
paired floors sum to $Q$, while all other pairs sum to $Q-1$.
The value at $s=0$ is $A$, establishing attainment and the formula.
\end{proof}

To apply the proposition to $S_{d,n}$, set $Q=d+n$. The affine map
\[
 x\longmapsto Q(x_1,\ldots,x_{d-1},0)
       +(1-\textstyle\sum_i x_i)\one_d
\]
sends $S_{d,n}$ to $Q\conv(0,e_1,\ldots,e_d)$ and its lattice to
$L=Q\Z^d+\Z\one_d$. The quotient $\Z^d/L$ identifies with
$(\Z/Q\Z)^{d-1}$ by subtracting the last coordinate, under which
the positive coordinate generators become those of $G_{d,Q}$.

For completeness, if $\Delta_d=\conv(0,e_1,\ldots,e_d)$ and
$L\subseteq\Z^d$ is a full-rank sublattice, write $x=b+f$ with
$b\in\Z^d$ and $0\le f_i<1$. The least $\rho$ for which
$x\in\rho\Delta_d+L$ equals
\[
 \sum_i f_i+
 \min\bigl\{\textstyle\sum_i u_i:
          u\in(b+L)\cap\Z_{\ge0}^d\bigr\}.
\]
The second term is the directed quotient distance of $b+L$.
Taking the supremum gives
$\mu(\Delta_d,L)=d+\operatorname{diam}(\Z^d/L)$, the standard
simplex/diameter identity of \cite[Lemma~B.1]{CSS}. Scaling by $Q$
and applying Proposition~\ref{d:graph-diameter} yields
\[
 \mu(S_{d,n})
 =\frac{\operatorname{diam}G_{d,Q}+d}{Q}
 =\frac{d-1}{2}+\frac{d+\gcd(d,n)}{2(d+n)}.
\]



\section{Exact-arithmetic certificates}
The following complete Python 3 code uses only the standard library.
For each ordered partition it constructs phases from the prescribed
waiting times and then evaluates the actual right-continuous fractional
flow at every crossing. Thus coincident jumps are checked using their
actual final values. The main text proves optimality for every input;
the finite checks here validate representative certificates and edge cases.
\begin{Verbatim}[fontsize=\footnotesize,frame=single]


"""Exact event certificates for the partition formula; Python 3 only."""
from fractions import Fraction as F
from itertools import product
from math import gcd


def frac(x):
    return x - x.numerator // x.denominator


def construct(k, weights, repeated, assignment):
    a = list(map(F, weights))
    b = F(repeated)
    A = sum(a)
    C = k*b + A
    B = (b+A)/2 - (k*b*b + sum(x*x for x in a))/(2*C)
    bins = [sorted([i for i, h in enumerate(assignment) if h == j],
                   key=lambda i: -a[i]) for j in range(k)]
    masses = [sum((a[i] for i in group), F(0)) for group in bins]
    D = [F(0)]
    for h in range(k):
        D.append(D[-1] + A/k - masses[h])
    assert D[-1] == 0
    Y = [v-min(D[:-1]) for v in D[:-1]]
    cost = b*sum(Y)
    phases = [F(0)]*(len(a)+1)
    for h, group in enumerate(bins):
        time, height, prefix = F(h, k), Y[h], F(0)
        for i in group:
            prefix += a[i]
            next_height = max(F(0), Y[h]-prefix)
            time += (a[i]+next_height-height)/C
            phases[i+1] = -time
            cost += a[i]*next_height
            height = next_height
        final_gap = (b+Y[(h+1) % k]-height)/C
        assert final_gap > 0
        assert time+final_gap == F(h+1, k)

    q, c = [k]+[1]*len(a), [b]+a
    events = {frac((F(j)-t)/r)
              for r, t in zip(q, phases) for j in range(r)}
    # This evaluates the actual right-continuous flow, including ties.
    actual = min(sum(w*frac(t+r*s)
                     for w, t, r in zip(c, phases, q)) for s in events)
    predicted = B-cost/C
    assert actual == predicted
    return predicted, B, len(events) < sum(q)


cases = [(2, [1, 2], 1), (2, [1, 1, 7], F(2, 3)),
         (3, [1, 2, 7], F(3, 2)),
         (4, [F(1, 3), F(2, 3), F(7, 4)], F(5, 6)),
         (5, [2, 3], 1), (1, [1, 2, 3], 2)]
count = ties = 0
for k, a, b in cases:
    best = None
    for assignment in product(range(k), repeat=len(a)):
        value, B, coincident = construct(k, a, b, assignment)
        best = value if best is None else max(best, value)
        count += 1
        ties += coincident
    if k == 2:
        a, b = list(map(F, a)), F(b)
        A, C = sum(a), 2*b+sum(a)
        delta = min(abs(A-2*sum((x for x, yes in zip(a, flags)
                               if yes), F(0)))
                    for flags in product([0, 1], repeat=len(a)))
        assert best == B-b*delta/(2*C)

for d in range(1, 6):
    for k in range(1, 6):
        a, b = F(2, 3), F(5, 4)
        assignment = tuple(h for h in range(k)
                           for _ in range((h+1)*d//k-h*d//k))
        value, _, _ = construct(k, [a]*d, b, assignment)
        expected = (d-1)*a/2 + a*b*(d+gcd(d, k))/(2*(k*b+d*a))
        assert value == expected

# Original smallest quotient and the nonintegral extension.
for r in map(F, [F(3, 2), 2, F(7, 3), 3, 10]):
    H = F(1, 2)+1/(r+1)
    conjectured = (r+2)/(2*r+1)
    assert H-conjectured == (r-1)/(2*(r+1)*(2*r+1))
    value, _, _ = construct(2, [1, 1], 1/r, (0, 1))
    assert value == H

for d in range(1, 13):
    for n in range(1, 31):
        h = (n+d-1)//d
        radius = F(d-1, 2)+F(d+gcd(d, n), 2*(d+n))
        target = F(d-1, 2)+F(1, h+1)
        assert radius <= target
        if h >= 2:
            assert (radius == target) == (n % d == 0)

print(f'Passed {count} exact constructed profiles, {ties} with ties.')
print('Weighted gcd, real-parameter gap, and interior-count checks passed.')


\end{Verbatim}


\begin{thebibliography}{9}
\interlinepenalty=10000
\bibitem{CSS}
G.~Codenotti, F.~Santos and M.~Schymura,
\emph{The covering radius and a discrete surface area for non-hollow
simplices}, Discrete \& Computational Geometry \textbf{67} (2022),
65--111. \href{https://doi.org/10.1007/s00454-021-00330-3}
{doi:10.1007/s00454-021-00330-3}.
\href{https://arxiv.org/abs/1903.02866}{arXiv:1903.02866}, v2 (2021).

\bibitem{Krivokuca}
K.~Krivoku\'ca, \emph{Upper bounds on covering minima of convex bodies},
\href{https://arxiv.org/abs/2601.15173}{arXiv:2601.15173v1},
21 January 2026.

\bibitem{GMS}
B.~Gonz\'alez Merino and M.~Henze,
\emph{On densities of lattice arrangements intersecting every
$i$-dimensional affine subspace},
\href{https://arxiv.org/abs/1605.00443}{arXiv:1605.00443v1} (2016).
References to Proposition~4.8 use this preprint version.

\bibitem{CGNS}
J.~Cslovjecsek, R.~D.~Malikiosis, M.~Nasz\'odi and M.~Schymura,
\emph{Computing the covering radius of a polytope with an application
to lonely runners}, Combinatorica \textbf{42} (2022), 463--490.
\href{https://doi.org/10.1007/s00493-020-4633-8}
{doi:10.1007/s00493-020-4633-8}.
\href{https://arxiv.org/abs/2009.12080}{arXiv:2009.12080v2}.

\bibitem{GJ}
M.~R.~Garey and D.~S.~Johnson,
\emph{Computers and Intractability: A Guide to the Theory of
NP-Completeness}, W.~H.~Freeman, 1979, problem SP15 (3-PARTITION).
See also their \emph{Complexity results for multiprocessor scheduling
under resource constraints}, SIAM Journal on Computing
\textbf{4} (1975), 397--411,
\href{https://doi.org/10.1137/0204035}{doi:10.1137/0204035}.

\bibitem{VMRC200}
\emph{Weighted fractional-part flows and sharp lattice covering bounds},
Verifiable Mathematical Reasoning Corpus, manuscript [200],
18 November 2026. Repository revision
\texttt{ac1492a50784be919fe0b5e8f19f96f8d221a22e}.
\href{https://github.com/QEEDDD/Verifiable-Mathematical-Reasoning-Corpus}
{Repository and manuscript source}.
\end{thebibliography}
\end{document}
